\BourbakiExercisesSection

\begin{enumerate}
\def\labelenumi{\arabic{enumi}.}
\tightlist
\item
  设 \(\Delta\) 是一个交换幺半群，用加法记号书写，其单位元记为 0。设 \(\mathbf{Z}^{(\Delta)}\) 是 \(\Delta\) 在 \(\mathbf{Z}\) 上的代数（III，§2，第 6 号），即一个交换环，它是自由 \(\mathbf{Z}\) -模，具有基 \((\mathrm{X}^{\sigma})_{\sigma \in \Delta}\) 和乘法 \((a\mathrm{X}^{\sigma})(b\mathrm{X}^{\tau}) = ab\mathrm{X}^{\sigma +\tau}\) 。对每个 \(\mathbf{Z}\) -模，我们记 \(\mathrm{E}^{(\Delta)} = \mathrm{E} \otimes_{\mathbf{Z}} \mathbf{Z}^{(\Delta)}\) ，因而 \(\mathrm{E}^{(\Delta)}\) 的每个元素都可唯一地写成形式 \(\sum_{\sigma \in \Delta} z_{\sigma} \otimes \mathrm{X}^{\sigma}\) ，其中 \(z_{\sigma} \in \mathrm{E}\) ，族 \((z_{\sigma})\) 有有限支撑。对每个 \(\mathbf{Z}\) -模同态 \(f: \mathrm{E} \to \mathrm{E}'\) ，用 \(f^{(\Delta)}\) 记 \(\mathbf{Z}^{(\Delta)}\) -模同态 \(f \otimes 1: \mathrm{E}^{(\Delta)} \to \mathrm{E}'^{(\Delta)}\) 。类似地，若 \(\Delta'\) 是另一个单位元记为 0 的加法幺半群，且 \(\alpha: \Delta \to \Delta'\) 是一个幺半群同态，则用 \(\alpha(\mathrm{E})\) 记同态 \(\mathrm{E}^{(\Delta)} \to \mathrm{E}^{(\Delta')}\) ，它满足 \(\alpha(\mathrm{E})(z_{\sigma} \otimes \mathrm{X}^{\sigma}) = z_{\sigma} \otimes \mathrm{X}^{\alpha(\sigma)}\) 。把 \((\mathrm{E}^{(\Delta)})^{(\Delta')}\) 与 \(\mathrm{E}^{(\Delta \times \Delta')}\) 等同，办法是把 \(\sum_{\sigma' \in \Delta'} \left( \sum_{\sigma \in \Delta} z_{\sigma, \sigma'} \otimes \mathrm{X}^{\sigma} \right) \otimes \mathrm{X}^{\sigma'}\) 对应到元素 \(\sum_{\sigma, \sigma'} z_{\sigma, \sigma'} \otimes \mathrm{X}^{(\sigma, \sigma')}.\) 最后，用 \(\varepsilon(\mathrm{E})\) 记同态 \(\mathrm{E}^{(\Delta)} \to \mathrm{E}\) ，它把 \(\sum_{\sigma} z_{\sigma} \otimes \mathrm{X}^{\sigma}\) 映到 \(\sum_{\sigma} z_{\sigma}\)
\end{enumerate}

\begin{enumerate}
\def\labelenumi{(\alph{enumi})}
\tightlist
\item
  设 E 是一个 Z-模， \((\mathrm{E}_{\sigma})_{\sigma\in\Delta}\) 是 E 上的一个 \(\Delta\) 型分次，且对所有 \(\sigma\in\Delta\) ，设 \(p_{\sigma}\) 是投影算子 \(E\to E_{\sigma}\) ，它对应于把 E 分解为诸 \(E_{\sigma}\) 的直和。对所有 \(x\in E\) ，设 \(\phi_{\mathbf{E}}(x)=\sum_{\sigma\in\Delta}p_{\sigma}(x)\otimes X^{\sigma}\) 。证明这样定义的同态 \(\phi_{E}\colon E\to E^{(\Delta)}\) 具有下列性质：
\end{enumerate}

\[
\varepsilon (\mathbf {E}) \circ \phi_ {\mathbf {E}} = 1 _ {\mathbf {E}}\tag{1}
\]

\[
\delta (\mathbf {E}) \circ \phi_ {\mathbf {E}} = \phi_ {\mathbf {E}} ^ {(\Delta)} \circ \phi_ {\mathbf {E}},\tag{2}
\]

其中 \(\delta: \Delta \to \Delta \times \Delta\) 是对角映射。

证明这样就定义了 E 上的 \(\Delta\) 型分次所成的集到 E 到 \(\mathrm{E}^{(\Delta)}\) 中的满足条件 (1) 和 (2) 的同态所成的集上的一个双射。

\begin{enumerate}
\def\labelenumi{(\alph{enumi})}
\setcounter{enumi}{1}
\item
  设 E、F 是两个 \(\Delta\) 型分次 Z-模。为使同态 \(f: E \to F\) 是 0 次分次的，必要且充分条件是 \(\phi_{F} \circ f = f^{(\Delta)} \circ \phi_{E}\) 。为使 E 的子模 \(E'\) 是分次的，必要且充分条件是 \(\phi_{E}(E') \subset E'^{(\Delta)}\) 。
\item
  设 \(\mathbf{A}\) 是一个环；给 \(\mathbf{A}^{(\Delta)} = \mathbf{A} \otimes_{\mathbf{Z}} \mathbf{Z}^{(\Delta)}\) 赋予一个环结构，使得 \((a \otimes \mathbf{X}^{\sigma})(b \otimes \mathbf{X}^{\tau}) = (ab) \otimes \mathbf{X}^{\sigma + \tau}\) 。设 \((\mathbf{A}_{\sigma})_{\sigma \in \Delta}\) 是加法群 A 的一个分次；对于这个分次与 A 上的环结构以及
\end{enumerate}

\(1 \in A_{0}\) 相容，必要且充分条件是 \(\phi_{A}: A \to A^{(\Delta)}\) 是一个环同态。陈述并证明关于分次 A-模的一个类似结果。

